% !TeX program = xelatex
\documentclass[11pt,a4paper]{article}
\usepackage[margin=23mm]{geometry}
\usepackage{fontspec}
\setmainfont{Latin Modern Roman}
\usepackage{amsmath,amssymb,mathtools}
\usepackage{xcolor,enumitem,booktabs,longtable,array,xurl,hyperref}
\definecolor{muted}{HTML}{53636C}
\hypersetup{colorlinks=true,urlcolor=blue,linkcolor=blue}
\setlength{\parindent}{0pt}
\setlength{\parskip}{5pt}
\newcommand{\R}{\mathbb R}
\newcommand{\E}{\mathbb E}
\newcommand{\N}{\mathbb N}
\newcommand{\ind}{\mathbf 1}
\newcommand{\Var}{\operatorname{Var}}
\newcommand{\Cov}{\operatorname{Cov}}
\newcommand{\supp}{\operatorname{supp}}
\DeclareMathOperator*{\argmin}{arg\,min}
\DeclareMathOperator*{\argmax}{arg\,max}
\newcommand{\entrymeta}[3]{{\sffamily\small\color{muted}\textbf{\texttt{#1}}\quad|\quad #2\par #3\par}}
\newcommand{\Problemref}[1]{\texttt{#1}}
\newcommand{\JELref}[2]{\texttt{#2} (JEL \texttt{#1})}
\begin{document}
\section*{Sequential randomization for mediation with intermediate covariates}
\textbf{Problem ID:} \texttt{C21-232}\quad \textbf{JEL:} \texttt{C21}\par
% BEGIN ECONOMICS STATEMENT
\label{problem:OP-0514}
\label{beyond:CI-35}
\entrymeta{CI-35}{Econometric theory and causal inference}{Research agenda formalization}
\subsection*{Definitions and assumptions}
First randomize binary $W$ conditional on baseline $X$, observe a post-treatment covariate $L$, and then randomize a feasible mediator setting $M$ conditional on $(X,W,L)$. Let $L(w)$ and $Y(w,l,m)$ be well-defined potential variables with consistency. A separate natural-mediator arm estimates a reference distribution $g_{w'}(m\mid X)$ of the mediator under treatment $w'$. Use a finite mediator state space, bounded outcomes, and known randomization probabilities $q(w\mid X)$ and $r(m\mid X,W,L)$ bounded away from zero. Let $X\in[0,1]^d$, $L\in[0,1]^{d_L}$. Require outcome means, $L$ densities, and reference mediator probabilities to have H\"older indices $\beta,\alpha,\gamma>0$, respectively, with radius $B$, and bounded intervention costs.

Write $\mathcal M$ for the finite mediator state space and $M^{\rm nat}(w)$ for the mediator under treatment $w$ with $L$ taking its natural value $L(w)$ and with no mediator intervention. The notation $M(w)$ below means $M^{\rm nat}(w)$, and
\[
 g_w(m\mid x)=\Pr\{M^{\rm nat}(w)=m\mid X=x\}
\]
is a population kernel, distinct from an estimator obtained from the reference arm. These potential variables and kernels refer to the target population; using a reference arm from another population requires justification that it recovers these kernels.

In the sequential arm, consistency means $L=L(W)$ and $Y=Y(W,L(W),M)$. Sequential randomization means that the treatment assignment is independent of the subject's potential variables conditional on $X$, and that the mediator assignment is independent of those potential variables conditional on $(X,W,L)$, using the stated assignment kernels. Define
\[
 \mu_w(x,l,m)=\E[Y\mid X=x,W=w,L=l,M=m],\qquad
 f_w(l\mid x)=\frac{dP\{L(w)\in\cdot\mid X=x\}}{dl}(l),
\]
where $dl$ is Lebesgue measure on $[0,1]^{d_L}$; conditional functions are understood almost everywhere under the relevant laws. The outcome means and $L$ densities mentioned above are $\mu_w$ and $f_w$.

Let $\mathcal F(x,w,l)\subseteq\mathcal M$ denote the feasible mediator settings. For each requested pair $(w,w')$, the displayed intervention and its positivity requirement concern every $m$ with $g_{w'}(m\mid x)>0$: such $m$ must belong to $\mathcal F(x,w,l)$ and have $r(m\mid x,w,l)>0$ for $P_X(dx)f_w(l\mid x)dl$-almost every $(x,l)$. If this support compatibility fails, the displayed target is not covered by the stated identification claim; conditioning or truncating the reference kernel would define a different target.
\subsection*{Formal question}
Define a stochastic mediator effect using a reference draw $G_{w'}\mid X\sim g_{w'}(\cdot\mid X)$ that is conditionally independent of the subject's potential variables given $X$:
\[
 \psi(w,w')=\E[Y(w,L(w),G_{w'})].
\]
Find cost-constrained sequential designs $(q,r)$ minimizing worst-case risk for contrasts among these means. Establish identification and characterize efficiency bounds and conditions for their attainment when $L$ both responds to $W$ and predicts the outcome, and characterize adaptive allocation rules retaining uniform inferential validity. For a prespecified real array $a=(a_{ww'})_{w,w'\in\{0,1\}}$ with $\sum_{w,w'}a_{ww'}=0$, write
\[
 \theta_a=\sum_{w,w'}a_{ww'}\psi(w,w'),\qquad
 \eta(w,w')=\E[Y(w,L(w),M^{\rm nat}(w'))],\qquad
 \eta_a=\sum_{w,w'}a_{ww'}\eta(w,w').
\]
The natural-effect comparison uses the same coefficient array: determine what extra structural restrictions, if any, make $\theta_a=\eta_a$ for the specified arrays. Here $L$ follows treatment $w$ in the outcome argument, while $M^{\rm nat}(w')$ includes its own natural intermediate covariate under $w'$; no other path-specific substitution is implicit. Equality of these mean contrasts is the requested comparison, not equality of distributions or an assumed cross-world independence condition.
\subsection*{Scope and status}
This precise design question is editorial. Randomization of the mediator setting is an intervention on a stated mechanism; an enhancement that merely shifts the mediator requires additional exclusion and compliance analysis. Sequential randomization identifies the displayed interventions under consistency and positivity. It does not impose independence between $Y(w,l,m)$ and a mediator counterfactual under another treatment, nor determine their cross-world coupling. Accordingly, natural effects such as $\E[Y(w,L(w),M(w'))]$ are not claimed identified merely from two randomized assignments. Post-treatment $L$ is propagated under its treatment-specific law rather than held fixed without explanation.
\subsection*{References for the source of the problem}
Carlos Cinelli, Avi Feller, Guido Imbens, Edward Kennedy, Sara Magliacane, and Jose Zubizarreta, \emph{Challenges in Statistics: A Dozen Challenges in Causality and Causal Inference}, arXiv:2508.17099v1 [stat.ME], 23 August 2025, \S3.4.3, ``Designing experiments for mediation analysis''; \S3.4.2, ``Identification \& estimation,'' for cross-world restrictions. \url{https://arxiv.org/html/2508.17099v1#S3.SS4.SSS3}.

\subsection*{Common conventions: Source mathematical conventions}

\textbf{Well-defined objects.}
All probability spaces and measurable variables are specified or inherited
from a local statistical model. Expectations used as finite targets require
the stated integrability. A decision rule or estimator is measurable with
respect to its available information; randomized rules may use an auxiliary
independent random variable. Norms without a subscript are Euclidean in
finite-dimensional spaces. An optimizer is requested only when attainment
is assured; otherwise the question concerns the optimal value and
$\varepsilon$-optimal admissible rules for every $\varepsilon>0$.

\textbf{Identification.}
A structural model $m\in\mathcal M$ implies an observable law
$P_m^{\rm obs}$ and target $\theta(m)$. The target is point identified on
$\mathcal M$ if
\[
 P_{m_1}^{\rm obs}=P_{m_2}^{\rm obs}
 \quad\Longrightarrow\quad\theta(m_1)=\theta(m_2)
 \qquad(m_1,m_2\in\mathcal M).
\]
When this fails, the sharp identified set is
\[
 \Theta(P;\mathcal M)=\{\theta(m):m\in\mathcal M,
                                  P_m^{\rm obs}=P\}.
\]
Specifying an intervention or estimand does not establish identification.
Positivity, exclusion, causal sufficiency, measurement restrictions, and
transport assumptions are substantive local inputs, never automatic
consequences of notation.

\textbf{Minimax risk and nontrivial inference.}
For a statistical experiment with data law $P^{(n)}$, target $\theta(P)$,
and nonnegative loss $\ell$, define
\[
 \mathcal R_n(\mathcal P,\ell)
   =\inf_{\widehat\theta_n}\sup_{P\in\mathcal P}
      \E_{P^{(n)}}\ell(\widehat\theta_n,\theta(P)).
\]
The infimum is over measurable estimators. A sharp rate is a sequence
$r_n$ for which $0<c\leq C<\infty$, independent of $n$, satisfies
$c r_n\leq\mathcal R_n\leq C r_n$ for all sufficiently large $n$.
Squared-error parametric risk is of order $n^{-1}$; the corresponding
root-mean-squared error is of order $n^{-1/2}$. These different conventions
are distinguished locally. A uniform asymptotic confidence region $C_n$
must satisfy
\[
 \liminf_{n\to\infty}\inf_{P\in\mathcal P}
     \Pr_{P^{(n)}}\{\theta(P)\in C_n\}\geq 1-\alpha,
 \qquad 0<\alpha<1,
\]
together with an informative diameter or expected-width requirement.
Coverage alone permits the vacuous whole parameter space. Testing questions
similarly impose both size control and power against a declared separated
alternative class.

\textbf{Binary causal baseline.}
When an entry explicitly invokes this baseline, one observes independent
copies of $O=(X,W,Y)$, with $W\in\{0,1\}$ and potential outcomes $Y(0),Y(1)$.
Assume consistency $Y=Y(W)$, no interference, conditional exchangeability
$(Y(0),Y(1))\perp W\mid X$, and overlap
$\epsilon\leq e(X)\leq1-\epsilon$ almost surely for a fixed
$0<\epsilon<1/2$, where
\[
 e(x)=\Pr(W=1\mid X=x),\qquad
 m_w(x)=\E[Y\mid W=w,X=x],\qquad
 \tau(x)=m_1(x)-m_0(x).
\]
These assumptions identify conditional mean effects almost everywhere
under the covariate law. Pointwise or uniform effects require appropriate
regularity and a specified support. Continuous treatment, longitudinal
data, adaptive experiments, and network interference require their own
assumptions in place of this baseline. Bounded outcomes, densities, and
smoothness are additional assumptions only where stated.

\textbf{Function classes.}
On $[0,1]^d$, $\mathcal H_d^s(L)$ is the isotropic Hölder ball of radius
$L>0$: put $k=\lceil s\rceil-1$ and $\gamma=s-k\in(0,1]$, bound derivatives
through order $k$, and require the order-$k$ derivatives to be
$\gamma$-Hölder with constant at most $L$. Thus integer orders use a
Lipschitz final derivative convention. The class
$\operatorname{BV}_d(L)$ consists of functions $f\in L^1((0,1)^d)$ whose
distributional derivative is a finite vector measure and for which
$\|f\|_{L^1}+|Df|((0,1)^d)\leq L$.
An $L^\infty$ bound or a particular pointwise version is imposed separately
when needed. Sparse and neural-network classes require a declared
dictionary or architecture and coefficient bounds.

An anisotropic H\"older class with declared block smoothness indices means
coordinatewise H\"older regularity in each block, uniformly in the other
blocks, with all corresponding derivative and H\"older bounds at most
the stated radius. Derivative orders and the integer-order convention in
each block are those just given. Mixed-derivative bounds are additional
restrictions only when explicitly stated.

\textbf{Decision and computational questions.}
Policy regret compares admissible feasible policies under the same law and
constraints. In computational entries, finite data and thresholds have
rational encodings; running time is measured in their total bit length.
Real-valued equilibrium search uses the explicitly declared algebraic or
FIXP encoding. An approximation ratio for maximization is written
$\operatorname{OPT}/\operatorname{ALG}$ and for minimization as
$\operatorname{ALG}/\operatorname{OPT}$, with zero optima and infeasible
instances handled locally. Historical approximation bounds are not
automatically current bounds.
% END ECONOMICS STATEMENT
\end{document}
